\section{Proof of the contraction theorem}
\label{app:tensor}

Theorem~\ref{thm:tensor-contraction} is proved by recursion on the
number of wires. Lemma~\ref{lem:tensor-arrays} balances the edge ends;
Lemma~\ref{lem:tensor-integrality} permits evaluation over the coefficient
ring; and Lemma~\ref{lem:decreasing-hermite-multiplicities} recovers each call's answer
from its local evaluations. The identity
\eqref{eq:tensor-recursion-identity} expresses each evaluation as a sum
of products of contractions with fewer than one eighth as many wires.
Section~\ref{subsec:tensor-cost} combines these facts and bounds the
complete recursion.

A recursive call receives a bipartite network with $M$ edges and
maximum degree at most the bound $\Delta$ of
Section~\ref{subsec:tensor-packing}, whose entries lie in
the coefficient ring modulo $2^p$ for a precision $p$, together with its depth $d$ in the
recursion; it returns the contraction modulo $2^p$. The root call
has $d=0$ and $p=P_0$. A call with $M\le H_{\mathrm{leaf}}$ edges, for the threshold
$H_{\mathrm{leaf}}$ of Section~\ref{subsec:tensor-packing}, is a \emph{leaf} and sums
directly over all Boolean edge assignments. Every other call is
\emph{nonleaf} and performs the following steps, in which the array
length $L$ and the number $u$ of retained two-adic levels are those of
\eqref{eq:tensor-parameters}, and the grid $S$ and the scaling
$\alpha$ are those of Section~\ref{subsec:tensor-deformation}.
\begin{enumerate}
\item Lift its entries using least nonnegative coordinate representatives
      (Section~\ref{subsec:tensor-arithmetic}).
\item Label the vertices and place the edges into arrays
      (Lemma~\ref{lem:tensor-arrays}).
\item Write the contraction as the value $\mathcal D(0)$ of the deformed
      contraction $\mathcal D(T)$ in \eqref{eq:tensor-deformed-contraction},
      with the polynomial matrix $W(T)$. Call $t\in S\setminus\{0\}$ a
      \emph{selected node} if $v_2(t)<u$; these are the nodes of positive
      multiplicity $q_t$ in Lemma~\ref{lem:decreasing-hermite-multiplicities}. At a
      selected node $t$, the matrix $W(t)$ joins only row and column
      edge ends of vertices whose labels from the previous step agree in
      coordinate $v_2(t)$; the $r$ classes of this coordinate are the
      groups.
\item For each selected node $t$ and each of $M(2L-1)+1$ interpolation
      values $\zeta$, compute $\mathcal D(t+\alpha\zeta)$ modulo $2^{p_t}$ by
      \eqref{eq:tensor-recursion-identity}: the entries of $W$ joining
      different groups are expanded in auxiliary variables
      (Section~\ref{subsec:tensor-auxiliary-variables}), and each surviving tuple of
      exponents of the auxiliary variables contributes the product of the contractions of
      $r$ child networks with at most $M/8$ edges each
      (Section~\ref{subsec:tensor-children}). The children are calls
      at depth $d+1$ with precision $p_t$.
\item For each selected node $t$, recover $\mathcal D(t+\alpha z)$ modulo
      $2^{p_t}$ from these values by Lagrange interpolation; then
      recover $\mathcal D(0)$ modulo $2^p$ by Lemma~\ref{lem:decreasing-hermite-multiplicities}.
\end{enumerate}
The recovery in the last step divides by powers of two, so the
children need the precision $p_t>p$. This is why each call first
fixes lifts of its entries, which it receives only modulo $2^p$.

\subsection{The coefficient ring and exact arithmetic}
\label{subsec:tensor-arithmetic}

The recursion recovers polynomials in an interpolation variable modulo
a power of two from their values at fewer than $s^4$ points
(Section~\ref{subsec:tensor-hermite}). Interpolation modulo $2^p$ needs
points whose pairwise differences are units, and modulo two the integers
supply only two such points. We therefore work with polynomials in a
root $U$ of a polynomial $f$ of degree $e$ that is irreducible modulo
two, with $2^e\ge s^4$.
Write $\lambda=\lceil\log_2 s\rceil$ and $e=4\lambda$.
Choose a monic polynomial $f\in\mathbb Z[U]$ of degree $e$, with
coefficients in $\{0,1\}$ below its leading coefficient, whose reduction
modulo two is irreducible.  This can be done deterministically in
polynomial time in $s$ by enumerating the $2^e$ candidates and testing
divisibility over $\mathbb F_2$ by every monic polynomial of positive
degree at most $e/2$. A reducible degree-$e$ polynomial has a factor
in this range. Both enumerations have $2^{O(e)}=s^{O(1)}$ terms.
Finite fields admit an irreducible polynomial of every positive
degree~\cite{vzGG13}; specifically, the number of monic irreducibles
of degree $e$ over $\mathbb F_2$ is at least
$2^e/(2e)>0$~\cite[Theorem~19.12]{Shoup09}.
Put
\[
 \mathcal O=\mathbb Z_{(2)}[U]/(f),\qquad
 \mathcal K=\mathbb Q[U]/(f),\qquad
 \mathcal A=\mathcal O[V]/(V^\ell),\qquad
 \mathcal A_p=\mathcal A/(2^p).
\]
Here $\mathbb Z_{(2)}$ consists of rational numbers with odd
denominator.  Irreducibility modulo two implies irreducibility over
$\mathbb Q$, so $\mathcal K$ is a field.  For a nonzero element of
$\mathcal K$, let $v_2$ be the minimum of the ordinary two-adic
valuations of its coordinates in $1,U,\ldots,U^{e-1}$.
Reduction of $\mathcal O$ modulo two is the field $\mathbb F_{2^e}$:
we have $\mathcal O/2\mathcal O\cong\mathbb F_2[U]/(\bar f)$, where
$\bar f$ is the reduction of $f$ modulo two, and this quotient is a
field with $2^e$ elements because $\bar f$ is irreducible of
degree~$e$. Now write $x=2^\alpha x_0$ and $y=2^\beta y_0$, where
$x_0,y_0\in\mathcal O$ have nonzero reductions modulo two.
Their product has nonzero reduction in this field, so
$v_2(x_0y_0)=0$. Thus
\[
 v_2(xy)=v_2(x)+v_2(y)\qquad(x,y\ne0).
\]
Every element of $\mathcal O$ that is nonzero modulo two is a unit because
the determinant of its multiplication matrix is invertible in
$\mathbb Z_{(2)}$, so the inverse matrix has entries in that ring.

Arithmetic in $\mathcal A_p$ takes polynomial time and space in
$s,p,\ell$.  We use the bases $1,U,\ldots,U^{e-1}$ and
$1,V,\ldots,V^{\ell-1}$, and represent each coordinate by its integer
residue between zero and $2^p-1$.  Inverting a scalar unit can be done
by inversion modulo two followed by Newton lifting, or by its
multiplication matrix, using exact polynomial and integer
operations~\cite{vzGG13}. For the lifting method, if $ab\equiv1$
modulo $2^j$, put $b'=b(2-ab)$. The identity
\[
 1-ab'=(1-ab)^2
\]
gives an inverse modulo $2^{2j}$ in the same coefficient ring.
Starting with the inverse in $\mathbb F_{2^e}$ and doubling the
precision therefore takes $O(\log(p+1))$ steps. This is the
prime-power inversion iteration of~\cite[Section~2.5, equation~(2.3)]{BZ10};
the displayed identity also proves it for $\mathcal O/(2^p)$.
Divisibility and precision for
elements of $\mathcal A$ are always coordinatewise.
We set $v_2(0)=+\infty$ and, for an element of $\mathcal K[V]/(V^\ell)$,
use the minimum valuation of its coefficients in $V$.

The binary-coordinate lifts of the $2^e$ residue-field elements give
interpolation points with pairwise unit differences.  Lagrange
interpolation at any required number of these points is therefore
valid over $\mathcal A_p$. Explicitly, for such points
$\zeta_1,\ldots,\zeta_N$ and $G\in\mathcal A_p[z]$ of degree less than $N$,
\[
 G(z)=\sum_{j=1}^N G(\zeta_j)
       \prod_{h\ne j}\frac{z-\zeta_h}{\zeta_j-\zeta_h}.
\]
Every denominator is a scalar unit. The evaluation matrix has
Vandermonde determinant $\prod_{j<h}(\zeta_h-\zeta_j)$, also a unit,
so the formula gives existence and uniqueness even with the
nilpotent variable $V$
\cite[Chapter~III, Exercise~1(3)--(5), p.~148]{LQ15}.
We will use fewer than $s^4$ points.
A contraction whose initial entries belong to $\mathcal B_{P_0,\ell}$
belongs to that subring after computation in $\mathcal A_{P_0}$,
because it is a polynomial in those entries.

When the precision increases, we lift each tensor entry separately.
For an entry originally specified modulo $2^p$, its lift in $\mathcal A$ uses
coordinate representatives in $\{0,\ldots,2^p-1\}$. This lifting need not preserve
addition or multiplication. Each recursive call contracts its specified
tensors modulo the precision of that call.  To define its deformation,
it first makes this entrywise choice of lifts.  All subsequent
evaluations in that call refer to these same lifts.
More explicitly, each tensor description is tagged with its creation
precision $p_c$ ($P_0$ for an original tensor).  At every later
precision, its entry is evaluated modulo $2^{p_c}$ and those same
integer representatives are used.  Precisions strictly increase
along a recursive path, so these representatives are also the
least nonnegative coordinate representatives at every later precision.  In particular,
an array tensor retains its creation-time representatives throughout the recursion.

The relation needed between calls is reduction, which is a ring
homomorphism. If $\rho_p:\mathcal A\to\mathcal A_p$ is reduction
and $Z$ denotes contraction, then for any chosen entrywise lifts
$\widetilde T_v$,
\[
 \rho_p\!\left(Z((\widetilde T_v)_v)\right)
       =Z((\rho_p(\widetilde T_v))_v).
\]
This follows term by term from the finite sum of products defining
$Z$. In particular, the deformation at the origin of a child call
reduces to the contraction specified by its parent, while its
evaluations at higher precision all use the child's fixed lifts.

\subsection{Parameters and packing into arrays}
\label{subsec:tensor-packing}

All estimates below hold for sufficiently large $s$, with the
threshold depending only on the fixed constants.  A fixed finite
range of smaller sizes is handled by direct summation. The parameter
inequalities can be tested directly.  Set
\[
 H_{\mathrm{leaf}}=\left\lceil\frac{s}{\lfloor\sqrt\lambda\rfloor}\right\rceil,
 \qquad r=64,\qquad \varepsilon=\frac1{128\cdot1024},
 \qquad \Delta=\lceil s^{9/10}\rceil.
\]
A call with $M\le H_{\mathrm{leaf}}$ edges is a leaf and is evaluated directly.
At a nonleaf call at depth $d$, with the root at depth zero,
$u_d$ is the number of retained valuation levels and $a_d$ is the
number of label coordinates, also used as slot bits. The cumulative
choice of $a_d$ will absorb inherited precision in
Section~\ref{subsec:tensor-hermite}. Set
\begin{equation}\label{eq:tensor-parameters}
 u=u_d=\left\lfloor\frac{\varepsilon\lambda}{2^d}\right\rfloor,
 \qquad a=a_d=2\sum_{j=0}^d u_j,\qquad
 a_{-1}=0,\qquad L=2^a,\qquad D_{\rm arr}=r^{2a}.
\end{equation}
We will arrange that every child has at most $M/8$ edges and maximum
degree at most $\Delta$.  These assertions imply, on every nonleaf path,
\[
 H_{\mathrm{leaf}}<M\le Cs\,8^{-d},\qquad
 d\le\frac{\log_2\lambda}{6}+O_C(1).
\]
Consequently
\begin{equation}\label{eq:tensor-parameter-bounds}
 u=\Theta_C(\lambda/2^d),\quad
 \lambda^{4/5}\le u\le M,\quad
 a\le4\varepsilon\lambda,\quad
 L,D_{\rm arr}\le s^{1/1000}
\end{equation}
for all sufficiently large $s$.  In particular, at each nonleaf call
\begin{equation}\label{eq:tensor-packing-conditions}
 M\ge2D_{\rm arr}^{\,2}\Delta,\qquad
 M\ge L D_{\rm arr},\qquad 2L\le\Delta.
\end{equation}
For example, the right side of the first inequality is
$O(s^{9/10+2/1000})=o(H_{\mathrm{leaf}})$.  The other inequalities follow in the
same way.

We call the two sides of the bipartite network its row side and its
column side, and speak of row and column vertices, edge ends and tensors
accordingly. In the deformation of Section~\ref{subsec:tensor-deformation}
the row edge ends index the rows of a matrix and the column edge ends its
columns.

The lemma below uses the following terms. Each vertex receives a label
in $[r]^a$, and the $i$th coordinate of its label is its \emph{symbol}
in coordinate $i$, for $0\le i<a$. An \emph{array} consists of $L$
\emph{slots}, numbered $0,\ldots,L-1$, and each slot has a row position
and a column position. An edge placed in a slot is attached at its row
end to the slot's row position and at its column end to the slot's
column position. The occupied positions are the edge ends;
unused slots carry no edge ends. Within an array, the original
identification of the two ends of each wire is thus the identity
between the row and column positions of the same slot. For a fixed
coordinate, the vertices on both sides with a given symbol in that
coordinate form a \emph{group}. At a nonzero interpolation node $t$,
Section~\ref{subsec:tensor-auxiliary-variables} uses the $r$ groups of coordinate
$v_2(t)$. The slot orientation in the second property below will make
the deformation local at that node; the edge-end bound in the third
property will give the child-size bound
(Section~\ref{subsec:tensor-children}).

\begin{lemma}[Packing into arrays]\label{lem:tensor-arrays}
For a nonleaf call with $M$ edges, maximum degree at most $\Delta$,
and the parameters in \eqref{eq:tensor-parameters} satisfying
\eqref{eq:tensor-packing-conditions}, the vertices on each side admit
labels in $[r]^a$ and the edges admit an embedding into arrays of
$L$ slots with the following properties.
\begin{enumerate}
\item There are at most $3M/L$ arrays.  Each edge occupies the same
      slot at its row and column ends. Unused slots have no edge ends.
\item For every array and every coordinate $i$, there are symbols
      $b\le c$ such that every edge in the array has the symbols $b$
      and $c$, in some order, at its row end and its column end in
      coordinate $i$. If $b<c$, the row end has symbol $b$ exactly
      when bit $i$ of the edge's slot number is zero. Consequently, a
      row edge end and a column edge end of the array belong to
      vertices with the same symbol in coordinate $i$ exactly when
      $b=c$ or their slot numbers differ in bit $i$.
\item For every coordinate and every symbol, at most $3M/r$ edge ends on each side belong to vertices with that symbol.
\end{enumerate}
A deterministic algorithm constructs these labels and this embedding in polynomial time and space in $s+\mathcal L$.
\end{lemma}
\begin{proof}
We give each vertex a single label in $[r]^a$, so all its edge ends
carry the same symbol in each coordinate. We balance the loads of
complete pairs of these labels, then pad them equally to arrange
the orientations by slot bits while keeping these vertex labels fixed.
For each pair of labels let $X_\gamma$ be the number of edges with
that pair.  Independent uniform vertex labels give
\[
 \mathbb E\sum_\gamma (X_\gamma-M/D_{\rm arr})^2\le2M\Delta.
\]
Indeed, expansion expresses this expectation as a sum over ordered
pairs of edges of their label-agreement probability minus
$1/D_{\rm arr}$.  Pairs with disjoint endpoints contribute zero.
At most $2M\Delta$ ordered pairs share an endpoint, including equal
edges, and each contributes at most one.
Conditional expectation gives a labeling with this bound.  To
implement it, try each of the $r^a$ labels of the next vertex and
compute the conditional expectation by enumerating pairs of edges.
Each probability concerns at most four vertices, so this is an
exact polynomial-time calculation.  This is the usual
conditional-expectation method~\cite{Raghavan88}.

Every load is at most $M/D_{\rm arr}+\sqrt{2M\Delta}$.
Pad all pair-label classes to a common number $H_{\rm pad}$ of
slots, the maximum load rounded up to a multiple of $L$.
The assumptions imply
\[
 H_{\rm pad}\le M/D_{\rm arr}+\sqrt{2M\Delta}+L
              \le3M/D_{\rm arr}.
\]
Fix the unordered pair of symbols in each coordinate.  If $a'$
coordinates have unequal symbols, their orientations give
$2^{a'}$ pair labels.  Partition their padded slots into arrays
using $2^{a-a'}$ slots from each of these labels per array.
This is possible because $H_{\rm pad}$ is divisible by $L$.
Number an array's positions from $0$ to $L-1$.  Its $i$th bit
specifies the orientation when the $i$th pair is unequal and the copy
when the pair is equal. For an unequal pair $\{b,c\}$ with $b<c$,
use row symbol $b$ and column symbol $c$ at slots whose $i$th bit
is zero, and the reverse symbols at slots whose $i$th bit is one.
Thus a row slot and a column slot have equal symbols in this
coordinate precisely when their $i$th bits differ.

There are $H_{\rm pad}D_{\rm arr}/L\le3M/L$ arrays.  In a fixed
coordinate, a given symbol occurs in exactly $1/r$ of all padded
pair labels on either side.  Its number of edge ends is therefore
at most $H_{\rm pad}D_{\rm arr}/r\le3M/r$.
\end{proof}

\subsection{A polynomial deformation and integral local coordinates}
\label{subsec:tensor-deformation}

Fix a nonleaf call, its arrays, and its entrywise lifted tensors.
Use
\[
 S=\{-L,-L+1,\ldots,L-1\},\qquad
 \mathcal C(T)=\prod_{t\in S}(T-t),\qquad \alpha=2^{a+1}.
\]
Within an array, for its row slot $x$ and column slot $y$, put
\begin{equation}\label{eq:tensor-polynomial-matrix}
 W_{xy}(T)=\frac{\mathcal C(T)}
                  {\mathcal C'(0)(T-(y-x))}.
\end{equation}
The factor in the denominator cancels a numerator factor, since
$y-x\in S$.  Thus this is a polynomial over $\mathbb Q$ of degree
$2L-1$.  Set entries between different arrays to zero and restrict
to the edge ends.  Under the original edge identification,
$W(0)=I$, since the diagonal entries are one and each other entry has
an uncancelled zero factor.

Let $A(Y)$ and $B(X)$ be the products of the lifted row and column
tensors, respectively, on independent edge-end assignments
$Y,X\in\{0,1\}^M$.  Write $l_0(z)=1-z$ and $l_1(z)=z$.  Define
\begin{equation}\label{eq:tensor-deformed-contraction}
 \mathcal D(T)=\sum_{X,Y\in\{0,1\}^M} A(Y)B(X)
                  \prod_{\text{row edge ends }x}l_{Y_x}((W(T)X)_x).
\end{equation}
Its coefficients lie in $\mathcal K[V]/(V^\ell)$, its degree is at
most $M(2L-1)$, and $\mathcal D(0)$ is the contraction of the lifted tensors.
Indeed, for Boolean $X,Y$ the product at $T=0$ is
$\prod_x l_{Y_x}(X_x)=\mathds{1}_{\{X=Y\}}$; the sum over $Y$
therefore restores the original edge identification.
In particular, reducing $\mathcal D(0)$ modulo the call's precision gives
the required answer.

\begin{lemma}\label{lem:tensor-integrality}
For every $t\in S$, the entries of $W(t+\alpha z)$ belong to
$\mathcal O[z]$, and $\mathcal D(t+\alpha z)$ belongs to $\mathcal A[z]$.
For every integer $p\ge1$ and $\zeta\in\mathcal O$ given by its
$e$ coordinates modulo $2^p$, each entry
$W_{xy}(t+\alpha\zeta)\bmod 2^p$ can be computed using polynomially
many bit operations in $s,p$.
\end{lemma}
\begin{proof}
The grid $S$ is a complete residue system modulo $2^{a+1}$.
For $t\ne q$ in $S$, $v_2(t-q)\le a$, and the counts of nonzero
differences of valuation $j$ are $2^{a-j}$ for $0\le j\le a$:
these are the residue classes divisible by $2^j$ but not by
$2^{j+1}$. The counts are independent of $t$. Hence
$v_2(\mathcal C'(t))=v_2(\mathcal C'(0))$.
In
\[
 \frac{\mathcal C(t+\alpha z)}{\mathcal C'(0)}
  =\frac{\alpha z\prod_{q\in S\setminus\{t\}}
                         (t-q+\alpha z)}{\mathcal C'(0)},
\]
each factor indexed by $q\ne t$ is coefficientwise divisible by
$2^{v_2(t-q)}$.  After canceling these powers against the
denominator, the displayed polynomial is divisible by $2^{a+1}$.
To obtain \eqref{eq:tensor-polynomial-matrix}, omit just the factor indexed by
$y-x$. If $y-x=t$, this removes $\alpha z$ and leaves an integral
polynomial. Otherwise $\alpha z$ remains and the omitted factor has
two-adic content $2^{v_2(t-y+x)}$ with $v_2(t-y+x)\le a$, so the
result is still divisible by two. Substitution in
\eqref{eq:tensor-deformed-contraction} proves the assertion for $\mathcal D$.

For the evaluation, first divide each factor
$t-q+\alpha z$ with $q\ne t$ by the known integer $2^{v_2(t-q)}$.
Both coefficients of the resulting linear polynomial are integers.
Cancel the remaining explicit powers of two against the denominator
as above, and invert its odd part modulo $2^p$. We can now evaluate
the product directly at the supplied residue of $\zeta$, including
$\zeta=0$. The input has $ep$ bits, and the product has $O(L)$
factors. All coordinate operations take polynomial time in $s,p$.
\end{proof}

\subsection{Hermite interpolation with decreasing multiplicities}
\label{subsec:tensor-hermite}

The next lemma isolates the precision estimate.  All power-series
truncations in its statement are ordinary truncations in $z$.

\begin{lemma}[Hermite interpolation with decreasing multiplicities]\label{lem:decreasing-hermite-multiplicities}
Let integers $u,M$ satisfy $1\le u\le a$ and $M\ge u$, and put
$h=\lceil3M/u\rceil$.
For $t\in S\setminus\{0\}$ set
\[
 q_t=h\max(0,u-v_2(t)).
\]
Suppose $\mathcal D\in(\mathcal K[V]/(V^\ell))[T]$ has degree at most
$M(2L-1)$, $\mathcal D(0)\in\mathcal A$, and
$\mathcal D(t+\alpha z)\in\mathcal A[z]$ for every $t\in S$.
For every integer $p\ge1$, the residue $\mathcal D(0)\bmod2^p$ is
determined by, and can be computed from, the data
\[
 \mathcal D(t+\alpha z)\pmod{(2^{p_t},z^{q_t})},
 \qquad
 p_t=p+8hL/2^{v_2(t)}+8M(a+1),\qquad q_t>0.
\]
Recovery takes polynomial time and space in $M,L,\ell,e$ and the
precisions.
\end{lemma}
\begin{proof}
We use the Hermite formula for evaluation at zero and bound its
denominators separately at each node. The decreasing multiplicities
control the loss from the other nodes; scaling by $\alpha$ makes
the local power-series division integral.
For a nonzero node $t$, write $i=v_2(t)$, $D_t=L/2^i$, and
$q(i)=h\max(0,u-i)$, so $q_t=q(i)$.
There are $L$ odd points in $S$, so
$\sum_{t\ne0}q_t\ge Lhu\ge3ML>M(2L-1)$.  Also
$h\le4M/u$, and hence $q_t\le4M$.
For each $t$ of positive multiplicity put, writing $t'$ for the other
nodes and $q_{t'}$ for their multiplicities,
\[
 E_t(T)=\prod_{t'\in S\setminus\{0,t\}}
                    \left(\frac{T-t'}{t-t'}\right)^{q_{t'}}.
\]
Hermite interpolation gives
\begin{equation}\label{eq:tensor-hermite-origin}
 \mathcal D(0)=\sum_{t:q_t>0}E_t(0)
       \left(\left[\frac{\mathcal D(t+\alpha z)}
                            {E_t(t+\alpha z)}\right]_{<q_t}\right)
                            (-t/\alpha).
\end{equation}
For completeness, replace zero in the right side by an indeterminate
$T$, and $-t/\alpha$ by $(T-t)/\alpha$.  Each summand has degree
less than $\sum q_t$.  At its own node it has the prescribed Taylor coefficients,
and at every other node it vanishes to the required order.
This verification may be made separately in each coefficient of
$V$, over the field $\mathcal K$. The difference from $\mathcal D$ is therefore divisible by
$\prod_{t\ne0}(T-t)^{q_t}$ and has smaller degree, so is zero.
This is the usual Chinese-remainder construction of Hermite
interpolation \cite[Section~5.6]{vzGG13}.

Each factor of $E_t(t+\alpha z)$ is a power of
$1+\alpha z/(t-t')$.  It is integral and has constant term one.
Its reciprocal truncated at $q_t$ is therefore integral, so
forming the quotient in \eqref{eq:tensor-hermite-origin} loses no
precision.

Fix $i=v_2(t)$ and $D=D_t$.  Directly from the product,
\begin{align}
 v_2(E_t(0))
 &=\sum_{t'\ne0,t}q_{t'}\bigl(v_2(t')-v_2(t-t')\bigr)\notag\\
 &=D\sum_{j=1}^{a-i}\frac{j}{2^j}
                           \bigl(q(i+j)-q(i)\bigr)
 \ \ge -6Dh. \label{eq:tensor-hermite-loss}
\end{align}
To justify the second equality, points of valuation less than $i$
contribute zero.  For $1\le j\le a-i$, there are $D/2^j$ points $t'$ with
$v_2(t')=i+j$, and the same number with $v_2(t-t')=i+j$.
The latter have $v_2(t')=i$, and these levels exhaust the points
other than $t$ with that valuation.  These counts follow from the
complete residue grid, whose nonzero integer differences have
valuation at most $a$.  Finally,
$q(i+j)-q(i)\ge-hj$ and $\sum_{j\ge1}j^2/2^j=6$.

Evaluation of a polynomial of degree below $q_t$ at $-t/\alpha$
loses at most $4M(a+1)$ further powers of two.  Thus replacing
each quotient coefficient by any lift of its residue modulo
$2^{p_t}$ changes its contribution in
\eqref{eq:tensor-hermite-origin} by an element of $2^p\mathcal A$.
More explicitly, an error $\eta\in2^{p_t}\mathcal A$ in coefficient
$z^j$, for $0\le j<q_t\le4M$, satisfies coordinatewise
\[
 v_2\!\left(E_t(0)(-t/\alpha)^j\eta\right)
 \ge p_t-6hD_t-j(a+1-i)\ge p.
\]
Here $i=v_2(t)$, so the last inequality uses both precision
increments in the statement. The inverse of $E_t(t+\alpha z)$
has integral coefficients and constant term one; multiplying by it
preserves the modulus of every supplied Taylor coefficient.
Summing these errors loses no further precision.  The reconstructed
sum is integral, since it differs from the integral value $\mathcal D(0)$
by an element of $2^p\mathcal A$.

One can perform this last combination by rational arithmetic in the
specified bases.  The multipliers $E_t(0)(-t/\alpha)^j$,
$j<q_t$, are products of polynomially many rational factors of
$O(a)$ bits, counting multiplicity.  Their numerator and denominator
lengths are therefore polynomial in $M,L,a$.  There are polynomially
many such coefficients, so even a product of their denominators is
an adequate polynomial-bit common denominator.  After addition,
cancel its powers of two and reduce the resulting integral
coordinates modulo $2^p$.  This proves the algorithmic assertion.
\end{proof}

At a call with input precision $p$, apply this lemma with the
parameters in \eqref{eq:tensor-parameters}.  For every selected $t$,
obtain the entire polynomial $\mathcal D(t+\alpha z)$ modulo $2^{p_t}$ by
evaluation at $M(2L-1)+1$ distinct binary-coordinate lifts of
residue-field elements and Lagrange interpolation.  Lemma~\ref{lem:tensor-integrality} permits this
interpolation over $\mathcal A_{p_t}$.  The number of points is
less than $s^4$ by \eqref{eq:tensor-parameter-bounds}.
The evaluation algorithm below passes precision $p_t$ to its
children.

For each selected node $t$, put $i=v_2(t)$ and $D_t=L/2^i$.
The input and child precisions satisfy
\begin{equation}\label{eq:tensor-precision-bound}
 p\le s\lambda^{O(1)}2^{a_{d-1}},
 \qquad
 p_t=O_C(MD_t/u).
\end{equation}
The first assertion holds at the root since $P_0\le Cs$ and
$a_{-1}=0$. A nonleaf call at depth $d'$ passes to its children the
precision $p_t=p+8hD_t+8M(a_{d'}+1)$. Here $h\le4M/u_{d'}\le4M$,
$M\le Cs$, $D_t\le L=2^{a_{d'}}$ and $a_{d'}\le4\varepsilon\lambda$,
so the increment $p_t-p$ is at most
$32Cs\,2^{a_{d'}}+8Cs(4\varepsilon\lambda+1)$. Summing the increments
along the path from the root, a call at depth $d$ has
\[
 p\le Cs+\sum_{d'=0}^{d-1}\Bigl(32Cs\,2^{a_{d'}}
        +8Cs(4\varepsilon\lambda+1)\Bigr).
\]
Its ancestors are nonleaf calls, so $a_{d'+1}-a_{d'}=2u_{d'+1}\ge2$ for
$d'+1<d$, and the terms $2^{a_{d'}}$ sum to at most
$\frac43\,2^{a_{d-1}}$. There are $d=O_C(\log\lambda)$ terms, and
$2^{a_{d-1}}\ge1$. For sufficiently large $s$, the fixed constant
$C$ is at most $\lambda$, so
\[
 p\le s\lambda^{O(1)}2^{a_{d-1}}.
\]
This proves the first assertion for every call, including leaves.
For the second, $q_t>0$ implies $i<u$ and hence
\[
 D_t\ge2^{a_d-u}=2^{a_{d-1}+u}.
\]
Since $M>H_{\mathrm{leaf}}$ and $u\ge\lambda^{4/5}$, this extra factor $2^u$
dominates every fixed power of $\lambda$.  Thus the first bound
on $p$, as well as $8M(a+1)$, is $O_C(MD_t/u)$.
The remaining increment $8hD_t$ has that bound directly.
All precisions remain polynomial in $s$ because
$a_d\le4\varepsilon\lambda$.

The decrease in multiplicities accounts for the factor $1/u$ in this estimate.
For comparison, assign the same multiplicity $q_*=\Theta(M)$ to
every level below $u$, and zero to the remaining levels.  At the
last retained level $i=u-1$, the identity
\eqref{eq:tensor-hermite-loss} would give the loss
\[
 -q_*D_t\sum_{j=1}^{a-i}\frac{j}{2^j}=-\Theta(MD_t).
\]
Our multiplicities instead change by only $h=\Theta(M/u)$ per
level, giving loss $O(MD_t/u)$.  The cumulative choice of $a_d$
provides the additional factor $2^u$ needed to absorb inherited
precision and the evaluation loss.  Thus the smaller bound holds
at every retained level and every recursion depth.

\subsection{Auxiliary variables for interactions between groups}
\label{subsec:tensor-auxiliary-variables}

Fix a selected $t$ and an interpolation value $\zeta\in\mathcal O$
for $z$.  In this subsection $T=t+\alpha\zeta$ and all calculations
are modulo $2^{p_t}$.  Put $i=v_2(t)<u$ and $D=L/2^i$.
Group vertices on both sides by their symbol in coordinate $i$.
Entries of $W(T)$ whose two edge ends have the same group are called
local and are kept unchanged.

This choice of coordinate makes $W(t)$ local.  Indeed,
$W_{xy}(t)=0$ unless $y-x=t$.  Adding an integer of valuation $i$
flips bit $i$, so a surviving pair has different $i$th bits.
When the array's two symbols in that coordinate differ, this
means that the row and column symbols agree. Equal pairs of symbols
already lie in one group.  Moving from $t$ to $t+\alpha\zeta$ introduces
the crossing terms that we now expand.

In an array whose two symbols in coordinate $i$ differ, a row
position $x$ and a column position $y$ are in different groups
exactly when their $i$th bits are equal.  Partition columns into
residue classes
\[
 y=y_0+2^{i+1}j,\qquad 0\le y_0<2^{i+1}.
\]
The column group is constant on a residue class.  For a row crossing to
that residue class,
\[
 v_2(t-y_0+x)\le i.
\]
Otherwise $y_0$ would agree with $x+2^i$ modulo $2^{i+1}$,
contrary to equality of their $i$th bits. Put
$d_x=T-y_0+x$ and $v=v_2(t-y_0+x)\le i$.
For every $\zeta\in\mathcal O$, adding $\alpha\zeta$ preserves this
valuation, since $v_2(\alpha\zeta)\ge a+1>i$.
Thus $d_x/2^v$ is a unit and
$v_2(2^{i+1}j/d_x)\ge1$.
A geometric expansion consequently
gives
\begin{equation}\label{eq:tensor-linear-form-expansion}
 W_{xy}(T)=
 \sum_{\substack{b\ge0\\b+1<p_t}}
 \frac{\mathcal C(T)}
      {\mathcal C'(0)(T-y_0+x)}
 \left(\frac{2^{i+1}}{T-y_0+x}\right)^b j^b
 \pmod {2^{p_t}}.
\end{equation}
Indeed $\mathcal C(T)/\mathcal C'(0)$ has valuation at least
$a+1$ by the integrality proof, and the coefficient of $j^b$
in the displayed sum has valuation at least
\[
 (a+1-v)+b(i+1-v)\ge b+1.
\]
For $R_x=2^{i+1}j/d_x$, use the finite identity
\[
 (1-R_x)^{-1}=\sum_{b=0}^{p_t-2}R_x^b+
                  \frac{R_x^{p_t-1}}{1-R_x}.
\]
The last denominator is a unit. Multiplying its remainder by
$\mathcal C(T)/(\mathcal C'(0)d_x)$ gives an element of
$2^{p_t}\mathcal O$. This proves the expansion also when
$\zeta=0$ or $j=0$. These expressions are computable by
factoring the known powers of two and using only unit inverses.

For each such array, residue class, and $b$ introduce an auxiliary variable ${\omega}$
of weight $w_{\omega}=b+1$, to be replaced by
\[
 L_{\omega}(X)=\sum_{\text{occupied column positions }y\text{ in the residue class}}j^bX_y.
\]
As usual $j^0=1$, including $j=0$.
Its source is the residue class's column group and its target is the other
group in that array's coordinate pair.  Formula
\eqref{eq:tensor-linear-form-expansion} expresses each row form as its
local part plus a linear combination of its auxiliary variables with that target.
The coefficient of ${\omega}$ is divisible by $2^{w_{\omega}}$.
There are at most
\begin{equation}\label{eq:tensor-variable-multiplicity}
 (3M/L)2^{i+1}=6M/D
\end{equation}
variables of any one weight.

In the product of the affine factors
$l_{Y_x}((W(T)X)_x)$, every coefficient of a monomial in the auxiliary variables
${\omega}^\beta$ is divisible by $2^{\sum_{\omega}w_{\omega}\beta_{\omega}}$.
Consequently only the exponent tuples
\begin{equation}\label{eq:tensor-exponent-budget}
 \sum_{\omega}w_{\omega}\beta_{\omega}<p_t
\end{equation}
need be considered.  The number of these tuples is
\begin{equation}\label{eq:tensor-partition-count}
 \exp\bigl(O(\sqrt{Mp_t/D})\bigr)
       =\exp\bigl(O_C(M/\sqrt u)\bigr).
\end{equation}
Here is the complete counting estimate.  Fix $\theta>0$. Every tuple
$\beta$ satisfying \eqref{eq:tensor-exponent-budget} has
$\exp(\theta(p_t-\sum_{\omega}w_{\omega}\beta_{\omega}))\ge1$. Summing this quantity over
all tuples of nonnegative integers gives the generating product, so
the number of tuples is at most
\[
 \sum_{\beta}\exp\Bigl(\theta\bigl(p_t-\sum_{\omega}w_{\omega}\beta_{\omega}\bigr)\Bigr)
 =e^{\theta p_t}\prod_{\omega}\bigl(1-e^{-\theta w_{\omega}}\bigr)^{-1}.
\]
Every weight $w_{\omega}$ is a positive integer, and by
\eqref{eq:tensor-variable-multiplicity} at most $6M/D$ auxiliary variables have
any one weight. The number of tuples is therefore at most
\[
 \exp\left(\theta p_t+\frac{6M}{D}
                   \sum_{j\ge1}-\log(1-\exp(-\theta j))\right).
\]
Expanding the logarithm and summing the geometric series gives
\[
 \sum_{j\ge1}-\log(1-\exp(-\theta j))
  =\sum_{h\ge1}\frac1{h(\exp(\theta h)-1)}
  \le\frac1\theta\sum_{h\ge1}\frac1{h^2}.
\]
Take $\theta=\sqrt{M/(Dp_t)}$ and use
\eqref{eq:tensor-precision-bound}.  All auxiliary variables and their
exponents have polynomial-size descriptions.  Tuples satisfying
\eqref{eq:tensor-exponent-budget} can be enumerated with polynomial
delay and space by ordering the variables and successively choosing
each exponent subject to the remaining budget.  Every feasible partial
choice has the all-zero continuation, so the number of visited
partial choices is at most the number of complete tuples times
the number of variables, up to an additive constant.

\subsection{The child networks and their descriptions}
\label{subsec:tensor-children}

For one tuple $\beta$ in \eqref{eq:tensor-exponent-budget} construct
a child network for each group $g\in[r]$.
Place all current vertices in group $g$, from both sides of the
current network, on one side of the child.  Retain their edge ends and
their chosen lifted entries, now reduced modulo $2^{p_t}$.
On the other side put one vertex for each array, with an edge
to every retained edge end of that array.  Its indices are precisely
the row bits $Y$ and column bits $X$ of group $g$ in the array.

To define this new array tensor, form the product of the row
factors $l_{Y_x}$ for group $g$ in the array, with their arguments
expressed in local column bits and auxiliary variables.  Extract the
coefficient with exponents $\beta_{\omega}$ for all auxiliary variables targeted at
$g$ in that array.  Multiply it by
\[
 \prod_{\substack{{\omega}\text{ in this array}\\
                  \text{source}({\omega})=g}}L_{\omega}(X)^{\beta_{\omega}}.
\]
Empty products equal one. If there are no target rows, the extracted
coefficient is zero whenever a target exponent is positive. If the
source residue class of ${\omega}$ has no occupied column position, then $L_{\omega}(X)=0$ and its
power contributes one for $\beta_{\omega}=0$ and zero otherwise.

Every auxiliary variable has one target and one source.  Coefficient extraction
therefore factors over the array-group row products, and the
substitution for each auxiliary variable depends only on bits in its source
group. To see the factorization explicitly, denote the row product
for array $\kappa$ and target group $g$ by $R_{\kappa,g}(Y_g,X_g;{\omega})$,
where $Y_g$ and $X_g$ are that group's row and column bits in the array.
Its formal variables are precisely the auxiliary variables of that array
targeted at $g$. These sets of formal variables are disjoint as
$(\kappa,g)$ varies. Consequently, before substituting the auxiliary variables,
\[
 [{\omega}^\beta]\prod_{\kappa,g}R_{\kappa,g}
 =\prod_{\kappa,g}
    [{\omega}^{\beta|_{\kappa,g}}]R_{\kappa,g},
\]
where $\beta|_{\kappa,g}$ is the restriction to those auxiliary variables with that target.
After substitution, assign each factor $L_{\omega}(X)^{\beta_{\omega}}$ to
its source group. All factors in a group then depend only on
that group's edge-end bits, including its original tensor entries.
Distributing the product in
\eqref{eq:tensor-deformed-contraction}, and then summing the
independent bits within groups, proves the exact identity
\begin{equation}\label{eq:tensor-recursion-identity}
 \mathcal D(t+\alpha\zeta)
   =\sum_{\beta:\,\sum w_{\omega}\beta_{\omega}<p_t}
                       \prod_{g=1}^r Z_{g,\beta}
                  \pmod {2^{p_t}},
\end{equation}
where $Z_{g,\beta}$ is the contraction of the child just defined.
The omitted monomials vanish by their coefficient divisibility.

Each edge of a child joins a retained edge end of a group-$g$ vertex,
from either side of the current network, to the vertex of that edge end's
array. The edges of the child therefore correspond to the edge ends
of group-$g$ vertices on the two sides. For the coordinate $i$ that
defines the groups, Lemma~\ref{lem:tensor-arrays}(3) bounds these
edge ends by $3M/r$ on each side, so a child has at most
$6M/r=3M/32<M/8$ edges.  Old vertex degrees are unchanged, and a
new array vertex has at most $2L\le\Delta$ edge ends.  Thus the
invariants used in \eqref{eq:tensor-parameter-bounds} hold
inductively.

All new tensors are given implicitly.  To evaluate an array
tensor entry, expand its product of at most $L$ affine row
factors.  Each factor has polynomially many terms in $s$,
including its local term.  Enumerating these choices, checking
the desired exponents, and summing the resulting terms takes
\[
 \exp(O(L\log s))
\]
time and polynomial space.  The factors $L_{\omega}(X)^{\beta_{\omega}}$ are
computed by ordinary repeated squaring in $\mathcal A_{p_t}$.
Because $L\le s^{1/1000}$, this entry cost is
$\exp(O(s^{9/10}))$.

A description of a new tensor consists of its array data, the
tuple $\beta$, the interpolation value, and the creation precision.
These data and the matrix entries within the array suffice to compute its entry directly.  A retained tensor
is evaluated from its original description, or from its own
creation-time description if it was created earlier in the
recursion.  Its residue at that precision is then lifted using least nonnegative coordinate representatives to the current precision.  The old descriptions are
immutable and shared by reference.  Thus no composition of
entry routines builds up with the depth, and all entry routines
continue to have the required time and space bounds.
In particular, a new array entry uses the assigned edge-end bits and
the matrix entries, while the tensors incident to those edge ends remain
separate vertices; evaluating the array never calls their evaluators.
Only polynomially many vertices and polynomially much new
description data are added along a path.

\subsection{Time and space bounds}
\label{subsec:tensor-cost}

\begin{proof}[Proof of Theorem~\ref{thm:tensor-contraction}]
Embed the original entries into $\mathcal A_{P_0}$ and perform
the recursive algorithm above.  At a nonleaf call, for every node
of positive multiplicity and every required residue-distinct
interpolation value, compute \eqref{eq:tensor-recursion-identity}
by its child contractions.  Interpolate $\mathcal D(t+\alpha z)$ and apply
Lemma~\ref{lem:decreasing-hermite-multiplicities}.  At a leaf, sum directly over all
Boolean edge assignments.  Lemma~\ref{lem:tensor-integrality},
the recursion identity, and the recovery lemma prove correctness
by induction from the leaves.

The number of interpolation nodes, values, and groups at a call
is polynomial in $s$, independently of $\mathcal L$.
By \eqref{eq:tensor-partition-count}, the number of children of
a nonleaf depth-$d$ call is at most
\[
 s^{O(1)}\exp\left(
       O_C\left(\frac{s}{\sqrt\lambda}
                    \left(\frac{\sqrt2}{8}\right)^d\right)\right).
\]
Indeed $M\le Cs\,8^{-d}$ and
$u_d=\Theta_C(\lambda/2^d)$.  Multiplying these bounds over the
$O_C(\log\lambda)$ depths, or bounding the number of nodes on
each level in this way, gives at most
\[
 \exp\left(O_C(s/\sqrt\lambda)
                   +O_C(\lambda\log\lambda)\right)
       =\exp(O_C(s/\sqrt\lambda))
\]
calls in total.  At a leaf, direct summation uses at most $2^{H_{\mathrm{leaf}}}$
assignments and polynomially many entry evaluations per assignment.
Its cost is bounded by a polynomial in $s+\mathcal L$ times
\[
 \exp(O_C(H_{\mathrm{leaf}})+O_{C,A}(s^{9/10}))
   =\exp(O_{C,A}(s/\sqrt\lambda)).
\]
Local arithmetic, packing, descriptions, and recovery use
polynomial time per generated term, apart from the already
bounded entry computations.  The description-dependent factor
bounds the work per call. Since the branching bound is independent of
$\mathcal L$, the total work retains a fixed polynomial factor in $s+\mathcal L$.
This proves the time bound.

For space, traverse the recursion depth first, accumulating one
sum over $\beta$ and one product over groups at a time.
Store the polynomially many Taylor data needed for a call's
recovery, its current interpolation and auxiliary-variable indices, and its
graph and new tensor descriptions.  Their sizes are polynomial
in $s+\mathcal L$, including the precisions by
\eqref{eq:tensor-precision-bound}.  There are only
$O_C(\log\lambda)$ stack levels.  The common-denominator
calculation in Lemma~\ref{lem:decreasing-hermite-multiplicities} also uses only
polynomial-bit integers. Each call accumulates its contraction terms
modulo its own fixed precision.
Original entry routines use polynomial space by hypothesis,
and the new routines have the same property.
This proves the space bound.  Finally extract the base
coordinates in $\mathcal A_{P_0}$, which give the answer in
$\mathcal B_{P_0,\ell}$.
\end{proof}
